\section{Volume-independent perturbative stability}

\begin{theorem}[T3: proved volume-independent stability]
\label{thm:T3}
Let $\Lambda$ range over finite subsets or tori of a fixed
$D$-dimensional bounded-degree lattice.  Let $H_{0,\Lambda}$ be the
DET16 Hamiltonian \eqref{eq:H0}, and let $V_\Lambda$ be an even
finite-range interaction satisfying Definition~\ref{def:pert}
uniformly in volume.  Then there exists $t_*>0$, depending only on
$D$, the interaction range, and the local interaction-norm bound, but
not on $\Lambda$, such that
\[
 \gap(H_{0,\Lambda}+tV_\Lambda)\geq\frac12
\tag{5.1}\label{eq:stablehalf}
\]
for every sufficiently large volume and $|t|\leq t_*$.  After
decreasing $t_*$ once to cover the finitely many smaller volumes, the
same statement holds for every finite $\Lambda$.
\end{theorem}

\begin{proof}
Every local term $q_{\phi,x}$ is positive and annihilates the product
state, so $H_0$ is frustration free.  Theorem~\ref{thm:T1} gives
global gap $\gamma=1$.  Equation \eqref{eq:exactLTQO} verifies the full
Local-TQO condition \eqref{eq:MZLTQO} with
\[
 \Delta_0(\ell)=0
\]
at every scale, not merely asymptotically.  Equation
\eqref{eq:localgap} verifies the Local-Gap condition with
$\gamma(r)=1$.  A finite-range interaction has a compactly supported,
hence rapidly decaying, function $f_0$.  The local on-cluster Hilbert
space is finite dimensional.  Finally, all perturbation terms are
even, so spacelike-separated local terms commute rather than
anticommute; the locality and Lieb--Robinson estimates used in
Theorem~\ref{thm:MZ} apply to this even CAR interaction.

Apply Theorem~\ref{thm:MZ} to the perturbation $tV$.  Its strength is
$J=|t|J_V$, so one may set $t_*=J_0/J_V$.  This gives
\eqref{eq:stablehalf} for $L\geq L_0$.  There are only finitely many
lattice shapes below $L_0$ up to the fixed local geometry.  Each has a
simple isolated ground eigenvalue and gap one at $t=0$; finite-matrix
eigenvalues are continuous in $t$.  Taking the minimum of their
positive continuity radii and the large-volume value gives one
positive threshold for all finite volumes.
\end{proof}

\section{Dynamical finite gauge links: Lie--Schwinger gap theorem}
\label{sec:dynlinks}

Theorem~\ref{thm:T3} applies to prescribed parallel transporters but does
not itself add fluctuating link Hilbert spaces.  We therefore audit the
finite-group dynamical-link extension first against the local
Schrieffer--Wolff construction of Bravyi, DiVincenzo, and
Loss~\cite{BDL2011}, and then against the gap theorem proved by
Fröhlich and Pizzo~\cite{FP2020}.

Let each oriented open-chain link carry the electric basis of
$\C[\mathbb Z_N]$, $N=2$ or $4$, and let $U_\ell$ shift its electric
flux.  On the unconstrained ambient tensor product put
\[
 H_{0,\mathrm{dyn}}
 =\sum_x(\id-P_{\phi,x})+h_E\sum_\ell E_\ell^2,\qquad
 V=\sum_{\ell=(x,x+1)}K_\ell ,
\tag{5.2}\label{eq:dynH0}
\]
where
\[
 K_\ell=\sum_{\alpha=1}^{m}
 \left(c^\dagger_{x,\alpha}U_\ell c_{x+1,\alpha}
 +c^\dagger_{x+1,\alpha}U_\ell^\dagger c_{x,\alpha}\right).
\tag{5.3}\label{eq:dynhop}
\]
Here $m=16$ for DET16 and $m=4$ in the numeric analogue.  Assigning each
outgoing link to its left matter cell makes \eqref{eq:dynH0} a sum of
one-cell terms and \eqref{eq:dynhop} a two-cell interaction.  Its product
ground projector is the tensor product of the DET projectors and the
zero-electric-flux link projectors, and its local gap is
\[
 \Delta=\min(1,h_E).
\tag{5.4}\label{eq:dyndelta}
\]
Moreover $\norm{K_\ell}=m$, so the BDL interaction strength obeys
\[
 J=\max_x\sum_{\ell\ni x}\norm{K_\ell}\leq2m
\tag{5.5}\label{eq:dynJ}
\]
uniformly in chain length.

\begin{proposition}[BDL hypothesis repair and remaining obstruction]
\label{prop:BDLaudit}
For finite $\mathbb Z_2$ or $\mathbb Z_4$ links, the Gauss-restricted
Hilbert space is not a tensor product and hence is not directly in the
setup of Section~4 of Bravyi--DiVincenzo--Loss.  The ambient-cell
formulation \eqref{eq:dynH0}--\eqref{eq:dynJ} repairs this locality
hypothesis, and gauge invariance makes every Gauss sector reducing.
Nevertheless the cited BDL result does not supply the requested
volume-independent spectral-gap theorem or a numerical threshold.
\end{proposition}

\begin{proof}
The local Gauss constraints correlate a matter cell with its incident
links, so the physical subspace does not factor into local Hilbert spaces.
On the ambient space, by contrast, the left-link cell assignment gives
exactly the tensor product, one-cell unperturbed Hamiltonian, product
ground space, and two-cell perturbation assumed in BDL Section~4.1.
The determinant changes charge by sixteen and the dressed hopping
transports one unit of link flux with the matter charge, so all terms
commute with every Gauss generator.  Restriction to the physical sector
therefore preserves the spectrum originating from its invariant
subspace.  Adding a Gauss penalty is unnecessary; it would produce
overlapping commuting local terms and would no longer have BDL's
one-site unperturbed form.

For a second-order truncation $n=2$, BDL Theorem~3 normalizes
$\norm{V}_1=1$ and obtains
\[
 \epsilon_c
 =\min\left\{\frac{\Delta\beta}{4n^2},
              \frac{\Delta}{32\alpha}\right\},
\tag{5.6}\label{eq:BDLeps}
\]
where $\alpha$ and $\beta$ arise from big-$O$ estimates in their
Lemma~4.2 and are not evaluated numerically.  Writing
$\epsilon=tJ$ and using \eqref{eq:dynJ} gives only the symbolic
conditions
\[
\begin{aligned}
 m=4:\quad
 |t|&<\Delta\min\left\{\frac{\beta}{128},
                  \frac1{256\alpha}\right\},\\
 m=16:\quad
 |t|&<\Delta\min\left\{\frac{\beta}{512},
                  \frac1{1024\alpha}\right\}.
\end{aligned}
\tag{5.7}\label{eq:BDLsymbolic}
\]
No concrete number follows without values of $\alpha,\beta$.
More decisively, BDL Theorem~3 controls the \emph{ground-state energy}
with an extensive truncation error
$O(|\Lambda||\epsilon|^{n+1})$; it does not control the low-sector
spectral gap.  The paragraph following their Lemma~4.1 explicitly states
that the lemma says nothing about that gap.  Consequently
\eqref{eq:BDLsymbolic} cannot be promoted to a uniform lower bound on
$\gap(H_{0,\mathrm{dyn}}+tV)$.
\end{proof}

\begin{theorem}[Fröhlich--Pizzo Lie--Schwinger theorem, full form used
here]
\label{thm:FP}
Let
$\Hh^{(L)}=\bigotimes_{j=1}^{L}\Hh_j$ with
$\Hh_j\simeq\C^M$, where $M<\infty$ is independent of $L$.
Suppose the identical one-site operators $H_j\geq0$ have a
one-dimensional kernel $\C\Omega_j$ and satisfy
\[
 H_j|_{\{\C\Omega_j\}^{\perp}}\geq\id.
\tag{5.8}\label{eq:FPlocal}
\]
Let
\[
 K_L(\tau)=\sum_{j=1}^{L}H_j+
 \tau\sum_{\substack{I_{k;i}\subset\{1,\ldots,L\}\\
                     k\leq\bar k}}V_{I_{k;i}},
\tag{5.9}\label{eq:FPK}
\]
where $\bar k<\infty$ is fixed independently of $L$, every
$V_{I_{k;i}}$ is self-adjoint and supported on its indicated interval,
and $\norm{V_{I_{k;i}}}\leq1$.  Then there is
$\tau_{\rm FP}>0$, independent of $L$, such that for every real
$|\tau|<\tau_{\rm FP}$ and every finite $L$, $K_L(\tau)$ has a unique
ground state and
\[
 \gap K_L(\tau)\geq\frac12.
\tag{5.10}\label{eq:FPgap}
\]
\end{theorem}

This is Theorem~3.5 of
J.~Fröhlich and A.~Pizzo,
\emph{Lie--Schwinger Block-Diagonalization and Gapped Quantum Chains},
Communications in Mathematical Physics \textbf{375} (2020),
2039--2069~\cite{FP2020}.  Its conclusion is a spectral-gap estimate,
not a ground-energy truncation estimate, and is uniform in chain
length.

\begin{proposition}[Lie--Schwinger hypothesis audit for finite gauge
links]
\label{prop:FPaudit}
For $N_g\in\{2,4\}$, $h_E>0$, $m\in\{4,16\}$, and constant $\phi$, the
open-chain Hamiltonian \eqref{eq:dynH0}--\eqref{eq:dynhop} satisfies
every hypothesis of Theorem~\ref{thm:FP} on an ambient tensor product
after adjoining one decoupled endpoint link.  With
\[
 \Delta=\min(1,h_E),\qquad
 \tau=\frac{tm}{\Delta},
\tag{5.11}\label{eq:FPscale}
\]
therefore
\[
 |t|<t_{\rm FP}^{(m)}(h_E)
 :=\frac{\Delta}{m}\tau_{\rm FP}
 \quad\Longrightarrow\quad
 \gap\!\left(
 H_{0,\mathrm{dyn}}+tV
 \right)\big|_{\Hh_{\rm phys}}
 \geq\frac{\Delta}{2},
\tag{5.12}\label{eq:FPwindow}
\]
and the physical Gauss sector has a unique ground state.  The bound is
uniform in the number of clusters.
\end{proposition}

\begin{proof}
The audit is hypothesis by hypothesis:
\begin{table}[htbp]
\centering
\small
\begin{tabularx}{\textwidth}{@{}p{0.25\textwidth}X p{0.10\textwidth}@{}}
\toprule
\textbf{FP hypothesis}&\textbf{DET chain realization}&\textbf{Result}\\
\midrule
Finite tensor-product cells&
$\Hh_x=\F(\C^m)\otimes\C[\mathbb Z_{N_g}]$, of dimension
$2^mN_g$ independent of chain length.&Pass\\
Identical one-site $H_x$&
Put the outgoing link in cell $x$ and append one decoupled final link,
so every cell carries
$q_{\phi,x}+h_EE_x^2$.&Pass\\
Rank-one local vacuum&
$P_{\phi,x}$ and the zero-electric-flux projector both have rank one;
their tensor product is the unique cell vacuum.&Pass\\
Normalized local gap&
The cell gap is $\Delta=\min(1,h_E)>0$; dividing $H_x$ by $\Delta$
gives \eqref{eq:FPlocal}.&Pass\\
Fixed finite interval range&
$K_{x,x+1}$ acts on two adjacent cells, hence $\bar k=1$ independently
of volume.&Pass\\
Ordinary tensor locality&
In an ordered Jordan--Wigner representation the adjacent strings
cancel; the even bilinear hopping, including its left-cell link
operator, acts on exactly two cells.&Pass\\
Bounded self-adjoint interaction&
$K_{x,x+1}=K_{x,x+1}^{\dagger}$ and
$\norm{K_{x,x+1}}=m$; thus
$V_{x,x+1}=K_{x,x+1}/m$ has norm one.&Pass\\
Real small coupling&
After factoring out $\Delta$, the FP coupling is precisely
$\tau=tm/\Delta$.&Pass\\
Gauss restriction&
This is not an FP premise: all terms commute with every Gauss
generator, and the ambient vacuum is physical because
$m=0\pmod {N_g}$ with zero link flux.&Pass\\
\bottomrule
\end{tabularx}
\caption{Fröhlich--Pizzo hypotheses audited against the finite-link
DET chain.}
\label{tab:FPaudit}
\end{table}
The appended endpoint link is decoupled and has a unique zero-flux
ground state.  The enlarged Hamiltonian is therefore the original
ambient Hamiltonian tensored with this one-link factor; a lower bound
on the enlarged gap is also a lower bound on the original ambient gap.

The normalized Hamiltonian is exactly \eqref{eq:FPK} with
$\bar k=1$, so Theorem~\ref{thm:FP} gives a unique ambient ground state
and normalized gap at least $1/2$.  Every Gauss generator commutes with
the Hamiltonian.  At $t=0$ the unique ambient vacuum is in the trivial
Gauss sector, since both $m=4$ and $m=16$ vanish modulo $2$ and $4$.
Uniqueness and the discrete gauge eigenvalues keep the continued ground
state in that sector.  The physical spectrum is a subset of the ambient
spectrum and contains its ground state, so restriction cannot reduce
the excitation gap.  Restoring the factor $\Delta$ proves
\eqref{eq:FPwindow}.
\end{proof}

\subsection{What is and is not explicit in the FP window}

The Lie--Schwinger proof contains one explicit series constant:
\[
 \frac{e^{8a}-8a-1}{a}+e^{8a}-1=1,\qquad
 a=0.023320199830777650\ldots .
\tag{5.14}\label{eq:FPa}
\]
Appendix~A of \cite{FP2020} uses the explicit convergence restriction
\[
 |\tau|<\frac a{16}
 =0.001457512489423603\ldots .
\]
Its local-gap estimate also has the explicit sufficient inequality
\[
 8|\tau|+16|\tau|
 \sum_{\ell=3}^{\infty}
 \frac{|\tau|^{(\ell-2)/3}}{\ell}\leq\frac12,
\tag{5.15}\label{eq:FPgapcondition}
\]
whose positive boundary is
$\tau_{\rm gap}=0.046907169193748369\ldots$.
However, the induction closing Theorem~3.4 of \cite{FP2020} introduces
unevaluated universal constants $C,C'$ and repeatedly requires
``sufficiently small'' $|\tau|$.  A scrupulous parameterization is
therefore
\[
 \tau_{\rm FP}
 =\min\left\{\frac a{16},\tau_{\rm gap},\tau_{\rm ind}\right\}>0,
\tag{5.16}\label{eq:FPparametric}
\]
where $\tau_{\rm ind}>0$ denotes those volume-independent but
numerically unevaluated induction restrictions.  Thus the two concrete
model families have the cited, parametric windows
\[
\begin{array}{c|cc}
\toprule
&m=4\ \text{analogue}&m=16\ \text{DET16}\\
\midrule
h_E=\tfrac12&|t|<\tau_{\rm FP}/8&|t|<\tau_{\rm FP}/32\\
h_E\geq1&|t|<\tau_{\rm FP}/4&|t|<\tau_{\rm FP}/16\\
\bottomrule
\end{array}
\tag{5.17}\label{eq:FPwindows}
\]
No numerical value of $t_{\rm FP}^{(m)}$ follows unless
$\tau_{\rm ind}$ is evaluated.

The deterministic companion
\[
 \texttt{experiments/tfpt-discovery/
 det16\_dynamical\_links\_probe.py}
\]
passes all $149$ checks.  It verifies exact Gauss invariance, commuting
unperturbed terms, rank-one local vacua, the norms
\eqref{eq:dynJ}, the FP range and normalization hypotheses, the two
explicit constants \eqref{eq:FPa}--\eqref{eq:FPgapcondition}, the
second-order effective Hamiltonian, and fourth-order finite-volume
residual scaling.  A long-range mutant supported on the two endpoints
has diameter $L-1$ and visibly fails the fixed-$\bar k$ premise.
The smallest scanned gap remains $0.728854416$ through $|t|=0.4$ on
two and three four-mode clusters.  At $h_E\geq1$, that scan point has
normalized coupling $\tau=1.6$ for $m=4$ and $\tau=6.4$ for $m=16$,
both far outside the cited window; the scan is evidence against a
finite-volume closing there, not an extension of
\eqref{eq:FPwindow}.  The finite $\mathbb Z_2/\mathbb Z_4$ chain item is
therefore cited-verified with the parametric FP window, not with an
invented numerical threshold.

\section{The proposed explicit count and why it is conditional}

\begin{proposition}[Elementary relative-form gap bound]
\label{prop:relative}
Let $P=|\Omega_{\Lambda,\phi}\rangle
\langle\Omega_{\Lambda,\phi}|$ and $Q=\id-P$.  Suppose a self-adjoint
perturbation $V$ satisfies
\[
 PVQ=QVP=0,\qquad PVP=v_0P,
\tag{6.1}\label{eq:block}
\]
and
\[
 \left|
 \ip{\psi}{(V-v_0\id)\psi}
 \right|
 \leq c\,\ip{\psi}{H_0\psi}
 \quad\text{for every }\psi\in Q\Hh.
\tag{6.2}\label{eq:formbound}
\]
Then the ground vector remains $|\Omega_{\Lambda,\phi}\rangle$, and
\[
 \gap(H_0+tV)\geq1-c|t|
 \geq1-2c|t|.
\tag{6.3}\label{eq:linearbound}
\]
In particular, the weaker displayed bound is positive for
$|t|<t_{0,\mathrm{rel}}:=1/(2c)$.
\end{proposition}

\begin{proof}
Condition \eqref{eq:block} makes $P\Hh$ and $Q\Hh$ reducing subspaces.
The energy on $P\Hh$ is $tv_0$.  For normalized $\psi\in Q\Hh$,
Theorem~\ref{thm:T1} gives
$\ip{\psi}{H_0\psi}\geq1$, and
\[
\begin{aligned}
 \ip{\psi}{(H_0+tV-tv_0\id)\psi}
 &\geq
 (1-c|t|)\ip{\psi}{H_0\psi}\\
 &\geq1-c|t|.
\end{aligned}
\]
The min--max principle proves the first inequality in
\eqref{eq:linearbound}; the second is a weakening for $c|t|\geq0$.
\end{proof}

